\chapter{Bevezetés}

A digitális oktatás napjainkban egyre nagyobb szerepet kap az egyetemi és vállalati képzésekben egyaránt. Az online tanulási platformok lehetővé teszik, hogy a tanulni vágyók időben és térben rugalmasan sajátítsák el az ismereteket, miközben az oktatók hatékonyan kezelhessék a kurzusanyagokat és nyomon követhessék a résztvevők előrehaladását.

Jelen dolgozat egy Java alapú e-learning webalkalmazás tervezését és megvalósítását mutatja be, amelyet az oktatási és képzési folyamatok modernizálása, valamint széleskörű digitális támogatása céljából fejlesztettem ki. Az alkalmazás neve ME Learning, és célja, hogy egyetlen, könnyen kezelhető felületen keresztül tegye lehetővé a kurzusok kezelését, a prezentációs anyagok megtekintését, valamint az interaktív kvízek kitöltését.

A fejlesztés során különös figyelmet fordítottam a szerepkör-alapú hozzáférés-vezérlésre: az adminisztrátorok, oktatók és hallgatók eltérő jogosultságokkal rendelkeznek a rendszerben. Az oktatók kurzusokat hozhatnak létre, PDF és PowerPoint prezentációkat tölthetnek fel, kvízeket szerkeszthetnek, és kezelhetik a beiratkozott felhasználók listáját. A hallgatók böngészhetik a nyilvános kurzusokat, beiratkozhatnak, megtekinthetik az oktatási anyagokat és kitölthetik a kvízeket.

A dolgozat szerkezete a klasszikus szoftverfejlesztési folyamatot követi: a koncepcióalkotást a tervezési döntések bemutatása követi, majd az implementáció részleteit és a tesztelés eredményeit tárgyalom. A fejlesztés Spring Boot keretrendszer segítségével, Thymeleaf sablon-motorral és H2 beágyazott adatbázissal készült.
